\documentclass[12pt,a4paper]{report}

\usepackage[utf8]{inputenc}
\usepackage[T1]{fontenc}
\usepackage[french]{babel}
\usepackage{amsmath, amssymb, amsfonts}
\usepackage{graphicx}
\usepackage{float}
\usepackage{geometry}
\usepackage{hyperref}
\usepackage{setspace}
\usepackage{enumitem}
\usepackage{booktabs}
\usepackage{algorithm2e}
\usepackage{newtxtext,newtxmath} % Times New Roman
\usepackage{titlesec}
\usepackage{tikz}
\usetikzlibrary{arrows.meta, positioning, shapes.geometric, shadows, fit}

\geometry{margin=2.5cm}
\onehalfspacing % Interligne 1,5

% --- Configuration des polices et tailles ---
% Texte normal : 12pt Times New Roman (défini par \documentclass[12pt])

% Titres des chapitres : Gras, Centré, 18pt
\titleformat{\chapter}[display]
  {\centering\normalfont\fontsize{18}{22}\selectfont\bfseries}
  {\chaptertitlename\ \thechapter}
  {10pt}
  {\fontsize{18}{22}\selectfont}

% Titres des sections (Paragraphes) : Gras, 18pt
\titleformat{\section}
  {\normalfont\fontsize{18}{22}\selectfont\bfseries}
  {\thesection}
  {1em}
  {}

% Titres des sous-sections (Sous-paragraphes) : Normal, 12pt
\titleformat{\subsection}
  {\normalfont\fontsize{12}{14}\selectfont\mdseries}
  {\thesubsection}
  {1em}
  {}

\hypersetup{
    colorlinks=true,
    linkcolor=black,
    filecolor=magenta,      
    urlcolor=cyan,
    citecolor=black,
}


\begin{document}

% --- Préliminaires ---
\pagenumbering{roman}
\chapter*{Dédicace}
\addcontentsline{toc}{chapter}{Dédicace}
\begin{center}
    \vspace*{1cm}
    \textit{À nos très chers parents, \\
    Aucun mot ne saurait exprimer notre gratitude pour Você amour inconditionnel, vos sacrifices et votre soutien indéfectible tout au long de nos études. Ce travail est le fruit de vos prières et de votre confiance.} \\
    \vspace{1cm}
    \textit{À nos frères et sœurs, \\
    Pour vos encouragements constants, votre présence et la joie que vous apportez dans nos vies.} \\
    \vspace{1cm}
    \textit{À nos enseignants et mentors, \\
    Qui nous ont transmis le savoir avec passion et nous ont guidés vers l'excellence académique et professionnelle.} \\
    \vspace{1cm}
    \textit{À tous nos amis et compagnons de route, \\
    Pour les moments de partage, d'entraide et de solidarité durant ce parcours universitaire.} \\
    \vspace{1cm}
    \textit{À tous ceux qui, d'une manière ou d'une autre, ont contribué à l'aboutissement de ce projet.} \\
    \vspace{1cm}
    \textbf{Nous vous dédions ce modeste travail.}
\end{center}
\newpage

\chapter*{Remerciements}
\addcontentsline{toc}{chapter}{Remerciements}
L'aboutissement de ce projet de fin d'études a été rendu possible grâce au soutien et à la collaboration de nombreuses personnes que nous tenons à remercier chaleureusement.

Nous tenons tout d'abord à exprimer notre profonde gratitude à notre encadrant professionnel, \texttt{[Nom de l'encadrant professionnel]}, pour sa confiance, sa disponibilité et ses précieux conseils tout au long de ce projet. Son expertise technique et sa vision stratégique ont été déterminantes dans l'orientation du Provisioning Hub, nous permettant de confronter la théorie aux réalités complexes du monde industriel.

Nous remercions également notre encadrant académique, \texttt{[Nom de l'encadrant académique]}, pour son suivi rigoureux, ses critiques constructives et son soutien méthodologique. Sa rigueur intellectuelle nous a permis de maintenir un haut niveau d'exigence et de structurer notre réflexion de manière académique.

Nos remerciements s'adressent également à l'ensemble des membres de l'équipe technique de \texttt{[Nom de l'entreprise]} pour leur accueil chaleureux, leur patience et leur collaboration fructueuse. Leurs retours d'expérience sur les systèmes legacy ont été essentiels pour identifier les points de douleur et concevoir une solution adaptée.

Enfin, nous exprimons notre reconnaissance aux membres du jury qui nous font l'honneur d'évaluer ce travail. Leurs remarques et suggestions seront accueillies avec intérêt pour enrichir la qualité finale de ce mémoire.
\newpage

\chapter*{Abstract}
\addcontentsline{toc}{chapter}{Abstract}
In the context of modernizing Human Resources data management, this project presents the design and implementation of the \textbf{Provisioning Hub}, a declarative platform designed to synchronize HR data across multiple target systems (LDAP, Jira, etc.). The legacy ecosystem, characterized by fragmented Java-based integrations, suffered from significant development bottlenecks, a lack of unified traceability, and operational rigidity. 

The Provisioning Hub transforms this "code-heavy" process into a configuration-driven workflow. By implementing a declarative pipeline (\textit{Fetch $\rightarrow$ Filter $\rightarrow$ Transform $\rightarrow$ Send}) and integrating a dynamic policy engine based on SpEL, the system eliminates the need for redeployments when adding new destinations or updating business rules. Key technical contributions include the \textbf{Hub 360} view for end-to-end traceability, a \textbf{Save-First} persistence pattern to ensure data resilience, and the use of a \textbf{Dead Letter Queue (DLQ)} for robust error handling. The solution is built on a modern tech stack including Java 21, Spring Boot 3.3, Angular 19, PostgreSQL 17, and RabbitMQ 3.13. The results demonstrate a significant reduction in integration time and a total recovery of traceability across all provisioning flows.

\textbf{Keywords:} HR Provisioning, Declarative Platform, SpEL, Resilience, Hub 360, Event-Driven Architecture.
\newpage

\chapter*{Résumé}
\addcontentsline{toc}{chapter}{Résumé}
Dans le cadre de la modernisation de la gestion des données de Ressources Humaines, ce projet présente la conception et la réalisation du \textbf{Provisioning Hub}, une plateforme déclarative destinée à synchroniser les données RH vers divers systèmes cibles (LDAP, Jira, etc.). L'écosystème legacy, caractérisé par des intégrations fragmentées en Java, souffrait de goulots d'étranglement majeurs lors du développement, d'un déficit de traçabilité unifiée et d'une forte rigidité opérationnelle.

Le Provisioning Hub transforme ce processus « centré sur le code » en un flux basé sur la configuration. En mettant en œuvre un pipeline déclaratif (\textit{Fetch $\rightarrow$ Filter $\rightarrow$ Transform $\rightarrow$ Send}) et en intégrant un moteur de règles dynamiques basé sur SpEL, le système élimine la nécessité de redéploiements lors de l'ajout de nouvelles destinations ou de la modification de règles métier. Les contributions techniques majeures incluent la vue \textbf{Hub 360} pour une traçabilité de bout en bout, un modèle de persistance \textbf{Save-First} pour garantir la résilience des données, ainsi que l'utilisation d'une \textbf{Dead Letter Queue (DLQ)} pour une gestion robuste des échecs. La solution s'appuie sur une pile technologique moderne comprenant Java 21, Spring Boot 3.3, Angular 19, PostgreSQL 17 et RabbitMQ 3.13. Les résultats démontrent une réduction significative des délais d'intégration et une traçabilité totale de l'ensemble des flux de provisionnement.

\textbf{Mots-clés :} Provisionnement RH, Plateforme Déclarative, SpEL, Résilience, Hub 360, Architecture Orientée Événements.
\newpage

\chapter*{Glossaire}
\addcontentsline{toc}{chapter}{Glossaire}
\begin{description}
    \item[LDAP] \textit{(Lightweight Directory Access Protocol)} : Protocole standard utilisé pour l'accès aux services d'annuaire.
    \item[SpEL] \textit{(Spring Expression Language)} : Langage d'expression puissant utilisé par le framework Spring pour manipuler des objets à l'exécution.
    \item[DLQ] \textit{(Dead Letter Queue)} : File d'attente spécifique où sont envoyés les messages qui n'ont pu être traités après plusieurs tentatives.
    \item[CQRS] \textit{(Command Query Responsibility Segregation)} : Pattern architectural séparant les opérations de lecture des opérations d'écriture.
    \item[Hexagonal Architecture] : Style d'architecture visant à isoler le cœur métier des détails techniques (infrastructure, API, base de données).
    \item[Save-First Pattern] : Stratégie de persistance consistant à enregistrer l'état d'un échange avant d'effectuer un appel réseau pour prévenir toute perte de données.
    \item[Provisionnement] : Processus de création, modification et suppression d'identités et d'accès dans des systèmes informatiques.
\end{description}
\newpage

\listoffigures
\addcontentsline{toc}{chapter}{\listfigurename}
\newpage

\listoftables
\addcontentsline{toc}{chapter}{\listtablename}
\newpage

\tableofcontents
\newpage


% --- Introduction Générale ---
\begin{center}
    {\fontsize{24}{28}\selectfont \textbf{Introduction Générale}}
\end{center}
\addcontentsline{toc}{chapter}{Introduction Générale}
\label{chap:intro_gen}
\vspace{1cm}

Dans un environnement d'entreprise comptant plus de 40 000 collaborateurs, la gestion des identités et des accès (\textit{Identity and Access Management} - IAM) ne représente pas seulement un défi technique, mais un enjeu stratégique de sécurité et de productivité. Le provisionnement, processus consistant à synchroniser les données RH vers des cibles techniques (LDAP, Jira, etc.), est le moteur de l'opérationnalité des utilisateurs.

C'est dans ce contexte que s'inscrit le projet \textbf{Provisioning Hub}, conçu comme l'orchestrateur central entre les sources de données RH et les systèmes cibles. L'objectif est de transformer un processus historiquement basé sur le développement d'intégrations spécifiques en une plateforme de configuration de workflows. En adoptant un paradigme déclaratif, le Provisioning Hub permet de définir dynamiquement les flux de données, assurant ainsi une agilité accrue et une gouvernance des données renforcée.

Le présent rapport détaille la transition d'un écosystème legacy fragmenté vers cette solution moderne. Il est structuré comme suit : le \textbf{Chapitre 1} analyse le contexte et la problématique ; le \textbf{Chapitre 2} expose la conception et l'architecture ; le \textbf{Chapitre 3} décrit la réalisation technique ; et le \textbf{Chapitre 4} présente la validation et les résultats.

% --- Corps du rapport ---
\pagenumbering{arabic}

\chapter{Analyse du Contexte et Problématique}
\label{chap:contexte}

L'objectif de ce chapitre est d'exposer de manière rigoureuse le contexte industriel et technologique du projet, d'analyser les défaillances de l'infrastructure héritée et de justifier le besoin d'une solution moderne et unifiée de gestion du provisionnement. Nous étudierons comment la dépendance vis-à-vis de scripts fragiles, le cloisonnement des données en silos et les risques liés à l'absence d'automatisation ont conduit au développement du Provisioning Hub.

\section{Introduction et Contexte Industriel}

Dans les grandes organisations contemporaines, la transformation digitale s'accompagne d'une multiplication exponentielle des applications métiers, des plateformes collaboratives et des référentiels techniques. L'accès sécurisé et rapide à ces outils est indispensable pour garantir l'efficacité des collaborateurs. Ce domaine est couvert par la gestion des identités et des accès (\textit{Identity and Access Management} ou IAM).

Pour une entreprise comptant plus de 40~000 collaborateurs, répartis sur plusieurs sites et sous différents statuts contractuels (employés internes, prestataires externes, stagiaires, partenaires), la gestion manuelle ou artisanale des accès est impossible. Le cycle de vie d'un collaborateur au sein du système d'information se résume à travers le modèle \textbf{JML} (\textit{Joiner, Mover, Leaver}) :
\begin{itemize}
    \item \textbf{Joiner (Entrée)} : Création des comptes techniques et attribution des privilèges minimums requis dès le premier jour de travail.
    \item \textbf{Mover (Mutation)} : Ajustement dynamique des droits en fonction des changements de département, de poste ou de site géographique.
    \item \textbf{Leaver (Départ)} : Révocation instantanée et exhaustive de l'ensemble des accès pour prévenir tout risque d'intrusion post-emploi.
\end{itemize}

Le projet \textbf{Provisioning Hub} a été initié pour servir de pivot central entre le référentiel des Ressources Humaines (RH), agissant comme source unique de vérité, et les applications consommatrices d'identités (cibles techniques) telles que les annuaires d'entreprise LDAP, l'outil Jira de gestion de projets, ou les services de messagerie.

---

\section{Étude de l'Existant : L'Écosystème Legacy}
\label{sec:existant}

Avant l'introduction du Provisioning Hub, la gestion et la synchronisation des données d'identités reposaient sur une infrastructure fragmentée, combinant des services synchrones vieillissants et des scripts d'administration locaux.

\subsection{Description des Composants Hérités}
L'ancienne architecture reposait principalement sur trois briques :
\begin{enumerate}
    \item \textbf{ms-hraccess} : Un microservice miroir chargé d'extraire périodiquement les données brutes issues du progiciel de gestion RH (HR Access). Ce service n'appliquait aucune règle de gestion ni de filtre et se contentait d'exposer les données collaborateur.
    \item \textbf{Services SOAP (MYI-HRA)} : Un ensemble de services web basés sur le protocole SOAP (\textit{Simple Object Access Protocol}) chargés de transporter et de formater les données d'identité pour les transmettre aux applications tierces.
    \item \textbf{Scripts de traitement par lots (.bat)} : Pour la création d'un nouveau collaborateur ou la suppression d'un collaborateur sortant, l'exploitation reposait sur l'exécution automatique ou planifiée de scripts Batch Windows (\texttt{.bat}). Ces scripts utilisaient des commandes système basiques pour interagir avec les bases de données et les annuaires. Cette approche par scripts Batch présentait des failles de sécurité et d'exploitation majeures :
\begin{itemize}
    \item \textbf{Sécurité défaillante} : Utilisation de mots de passe temporaires définis de manière statique au sein des scripts et absence de chiffrement des flux de communication.
    \item \textbf{Absence de centralisation des erreurs} : Les incidents de synchronisation étaient enregistrés dans des fichiers logs locaux (\texttt{.log}) dispersés sur les serveurs d'exécution, empêchant toute visibilité transverse pour les équipes d'audit et de sécurité.
    \item \textbf{Instabilité et absence de mécanismes de reprise} : La moindre interruption réseau ou de base de données bloquait l'exécution des scripts de manière silencieuse, sans alerter les administrateurs et sans possibilité de reprise automatique (rollback).
\end{itemize}
\end{enumerate}

\subsection{Limites Techniques et Fragmentation du Flux}
Cette architecture hybride présentait des lacunes conceptuelles majeures :
\begin{itemize}
    \item \textbf{Fragibilité opérationnelle} : Les scripts \texttt{.bat} manquaient de mécanismes avancés de gestion des erreurs (pas de gestion des transactions SQL ni de rollback).
    \item \textbf{Couplage temporel fort et pannes en cascade} : La communication via SOAP imposait des flux synchrones. Si un système cible (comme Jira) était momentanément indisponible, la requête SOAP échouait, bloquant le traitement de tout le lot d'utilisateurs en cours de synchronisation.
    \item \textbf{Lourdeur du format XML} : Le protocole SOAP repose sur des enveloppes XML verbeuses, entraînant une surconsommation de bande passante et des temps de sérialisation élevés.
\end{itemize}

Le schéma de ce flux d'intégration synchrone hérité et ses faiblesses structurelles sont modélisés dans la Figure~\ref{fig:flux_legacy}.

\begin{figure}[H]
    \centering
    \begin{tikzpicture}[
        node distance=1.5cm and 2cm,
        block/.style={draw, fill=blue!10, rectangle, minimum height=1cm, minimum width=2.5cm, align=center, rounded corners, drop shadow},
        sync/.style={draw=red, <->, >=Stealth, thick, dashed, line width=1.2pt},
        async/.style={draw=green!60!black, ->, >=Stealth, thick, line width=1.2pt}
    ]
        % Nodes
        \node[block] (source) {ms-hraccess \\ (Source RH)};
        \node[block, right=of source] (soap) {Services SOAP \\ (MYI-HRA) \& \\ Scripts .bat};
        \node[block, right=of soap] (target) {LDAP / Jira \\ (Cibles Techniques)};
        
        % Arrows
        \draw[async] (source) -- node[above, font=\small] {Extraction} (soap);
        \draw[sync] (soap) -- node[above=0.1cm, text=red, font=\small] {Synchrone / Scripts} node[below=0.1cm, text=red, font=\small] {(Point de Blocage)} (target);
        
        % Annotation for coupling
        \node[draw=red, fill=red!10, dashed, below=0.6cm of soap, align=center, font=\small, rounded corners] (warning) {Point de défaillance unique :\\ Si la cible ou le script échoue,\\ les données sont perdues et non tracées};
        \draw[red, dashed, ->] (warning) -- (soap);
    \end{tikzpicture}
    \caption{Flux d'intégration legacy synchrone et exécutions manuelles}
    \label{fig:flux_legacy}
\end{figure}

---

\section{Problématique Opérationnelle et Analyse des Risques}
\label{sec:problematique}

L'exploitation à grande échelle de ce système hérité a généré une dette technique et des risques critiques pour l'organisation.

\subsection{Le Risque de Sécurité Majeur : Les Comptes Orphelins}
Le dysfonctionnement le plus critique concernait la phase de départ des collaborateurs (\textit{Offboarding}). Bien que la date de sortie soit saisie correctement par les Ressources Humaines dans l'application source, l'échec d'un script Batch ou une exception lors d'un appel SOAP empêchait fréquemment la suppression du compte dans l'annuaire LDAP ou dans l'Active Directory. 

Ces comptes d'anciens collaborateurs restés actifs, appelés \textbf{comptes orphelins}, constituent des vulnérabilités de sécurité majeures. Ils représentent un vecteur privilégié pour des intrusions externes non détectées, violant les règles de base de la norme de sécurité **ISO/IEC 27001** ainsi que le Règlement Général sur la Protection des Données (**RGPD**) concernant le contrôle strict des accès aux données personnelles.

\subsection{Le Cloisonnement des Données (Absence de Vue 360)}
Dans l'ancienne infrastructure, il n'existait aucune base de données centrale unifiée pour corréler les identités. Chaque système d'information cible gérait ses propres données de façon cloisonnée, formant des **silos de données**. 

Par conséquent, il était impossible d'obtenir une vision transverse d'un collaborateur (par exemple, connaître en temps réel l'ensemble des comptes ouverts sur LDAP, Jira et la messagerie pour un identifiant donné). En cas de contrôle ou d'audit de sécurité, les administrateurs devaient exécuter des requêtes manuelles sur chaque application cible et corréler les données sur des fichiers tableurs, un processus chronophage et propice aux erreurs.

\subsection{Rigidité Fonctionnelle et Coûts d'Exploitation}
\begin{itemize}
    \item \textbf{Goulot d'étranglement du développement} : Les règles d'habilitation et de transformation de données étaient codées directement en Java. Toute modification de règle imposait une modification de code, un déploiement de version et un redémarrage du service.
    \item \textbf{Opacité et boîte noire} : Les logs étaient dispersés sur plusieurs serveurs. Le diagnostic d'un compte non créé prenait plusieurs heures de recherche manuelle.
    \item \textbf{Coût financier lié aux licences} : Les comptes Jira et SaaS orphelins restaient comptabilisés comme actifs par les éditeurs de logiciels, entraînant des dépenses de licences injustifiées sur des comptes inactifs.
\end{itemize}

Pour résumer la gravité de ces failles, le Tableau~\ref{tab:analyse_risques_legacy} présente une cartographie des risques du système existant.

\begin{table}[H]
    \centering
    \caption{Analyse des risques de sécurité et financiers de l'infrastructure legacy}
    \label{tab:analyse_risques_legacy}
    \begin{tabular}{lp{5.5cm}p{5.5cm}}
        \toprule
        \textbf{Catégorie} & \textbf{Risque Identifié} & \textbf{Impact Métier / Technique} \\
        \midrule
        Sécurité & \textbf{Comptes orphelins actifs} (Leavers non révoqués dans l'AD/LDAP) & Accès illicite aux ressources internes après le départ de collaborateurs, non-conformité RGPD et ISO 27001. \\
        Financier & \textbf{Surfacturation de licences} (comptes Jira/SaaS non désactivés) & Coûts récurrents élevés pour des comptes inactifs non résiliés. \\
        Exploitation & \textbf{Fragilité des scripts .bat} (échecs de création/suppression non supervisés) & Erreurs d'intégration silencieuses (Joiners non ajoutés ou Leavers non supprimés). \\
        Gouvernance & \textbf{Systèmes cloisonnés en silos} (absence de vue consolidée à 360 degrés) & Incohérences d'accès indétectables d'une cible à l'autre, audits de conformité longs et manuels. \\
        Audit & \textbf{Dispersion des logs} sur plus de 10 serveurs de bases de données & Diagnostic d'incident laborieux (MTTR élevé), manque d'historique de conformité unifiée. \\
        \bottomrule
    \end{tabular}
\end{table}

---

\section{Avantages de la Solution Proposée : Le Provisioning Hub}
\label{sec:objectifs}

Le Provisioning Hub a été pensé pour remplacer cette architecture fragile par une plateforme moderne, robuste et industrialisée, offrant des avantages techniques et métiers significatifs.

\subsection{Automatisation Complète du Cycle JML}
Le Provisioning Hub prend en charge l'ensemble du cycle de vie du collaborateur de manière automatisée. Dès qu'un changement d'état (entrée, mutation, sortie) est détecté dans le système RH, un message événementiel est publié et traité immédiatement. La révocation des accès en cas de départ est garantie à 100\,\% sans intervention humaine.

Le cycle JML standardisé et son automatisation sont illustrés par la Figure~\ref{fig:jml_lifecycle}.

\begin{figure}[H]
    \centering
    \begin{tikzpicture}[
        node distance=1.5cm and 1.8cm,
        state/.style={draw, fill=green!10, circle, minimum size=2.2cm, align=center, drop shadow, font=\small},
        event/.style={font=\small\itshape, text=blue},
        arrow/.style={->, >=Stealth, thick, line width=1.2pt}
    ]
        % States
        \node[state, fill=blue!10] (joiner) {Joiner \\ (Entrée)};
        \node[state, right=of joiner, fill=orange!10] (mover) {Mover \\ (Mutation)};
        \node[state, right=of mover, fill=red!10] (leaver) {Leaver \\ (Départ)};
        
        % Transitions
        \draw[arrow] (joiner) -- node[above, event] {Création comptes} node[below, event] {(AD, Jira)} (mover);
        \draw[arrow] (mover) -- node[above, event] {Mise à jour} node[below, event] {profil / rôle} (leaver);
        
        % Self loops or other transitions
        \path (mover) edge [loop above, arrow] node[event] {Changement UO} (mover);
        
        % Final action
        \node[draw=red, fill=red!20, rectangle, below=0.8cm of leaver, align=center, rounded corners, font=\small, drop shadow] (revoke) {Révocation automatique\\ des accès (LDAP/SaaS)};
        \draw[arrow, red, -{Stealth[scale=1.2]}] (leaver) -- node[right, text=red, font=\small, align=left] {Détection de la\\ date de sortie} (revoke);
    \end{tikzpicture}
    \caption{Cycle de vie des collaborateurs (JML) et révocation automatisée}
    \label{fig:jml_lifecycle}
\end{figure}

\subsection{Fiabilité Comparée des Modèles de Communication}
Sur le plan académique, la différence de fiabilité entre l'ancien système (synchrone et en série) et le Provisioning Hub (asynchrone et découplé) peut être modélisée mathématiquement.

Soit un flux de provisionnement devant synchroniser une identité sur $N$ cibles techniques indépendantes. On note $p_i$ la probabilité de disponibilité opérationnelle de la cible $i$ au moment du flux.
Dans le modèle d'intégration synchrone en série hérité, le système global est considéré comme défaillant si une seule cible échoue. La fiabilité globale $R_{synchrone}$ est donc le produit des fiabilités individuelles :
\begin{equation}
R_{synchrone} = \prod_{i=1}^{N} p_i
\end{equation}

Si nous avons $N = 5$ cibles ayant chacune une fiabilité de $p_i = 0,98$ (ce qui représente 2\,\% de taux d'indisponibilité ou micro-coupure), la fiabilité globale du flux synchrone tombe à :
\begin{equation}
R_{synchrone} = (0,98)^5 \approx 0,904 \quad (90,4\,\%)
\end{equation}
Cela signifie qu'environ 10\,\% des flux d'intégration globale se terminent en échec ou en perte de données chaque jour.

Le Provisioning Hub résout cette faiblesse grâce à un modèle asynchrone découplé par des files d'attente (RabbitMQ). La fiabilité globale de la synchronisation vers chaque cible $i$ est isolée de la disponibilité des autres cibles, et consolidée par un mécanisme de retry automatique. La fiabilité de chaque branche d'intégration $i$ devient alors :
\begin{equation}
R_{asynchrone, i} = 1 - (1 - p_i)^k
\end{equation}
où $k$ représente le nombre de tentatives automatiques paramétrées sur la file d'attente. Avec $p_i = 0,98$ et seulement $k = 3$ tentatives, la fiabilité par cible monte à :
\begin{equation}
R_{asynchrone, i} = 1 - (1 - 0,98)^3 = 1 - 0,000008 = 0,999992 \quad (99,9992\,\%)
\end{equation}
Le système garantit ainsi la livraison éventuelle des données à toutes les cibles sans jamais bloquer le pipeline global de l'entreprise.

\subsection{Hybridation des Flux : Workflow Synchrone Manuel et Asynchrone Événementiel}
Le Provisioning Hub résout l'opposition traditionnelle entre les flux en temps réel et la robustesse de l'arrière-plan en implémentant une double dynamique d'exécution :
\begin{itemize}
    \item \textbf{Workflow asynchrone événementiel (mode automatique)} : C'est le comportement nominal pour le traitement en masse et l'automatisation. Lors de l'ingestion d'un événement provenant de l'une des 7 sources (comme un changement RH ou une affectation de planning), le flux est traité de manière asynchrone. L'Orchestrateur persiste l'événement, évalue les politiques métier, crée le Job et publie un message dans RabbitMQ. Des agents d'exécution (\textit{workers}) dépilent ces messages de manière distribuée. Cela protège le système contre les surcharges réseau et assure une livraison garantie sans couplage temporel (Figure~\ref{fig:flow_async_hub}).
    
    \begin{figure}[H]
        \centering
        \begin{tikzpicture}[
            node distance=1.2cm and 1.8cm,
            block/.style={draw, fill=blue!10, rectangle, minimum height=0.8cm, minimum width=2.2cm, align=center, rounded corners, drop shadow, font=\small},
            db/.style={draw, fill=orange!10, cylinder, minimum height=1.2cm, minimum width=1.4cm, align=center, shape border rotate=90, drop shadow, font=\small},
            queue/.style={draw, fill=green!10, rectangle, minimum height=0.8cm, minimum width=1.8cm, align=center, rounded corners, drop shadow, font=\small},
            arrow/.style={->, >=Stealth, thick, line width=1.2pt}
        ]
            % Nodes
            \node[block] (source) {Source Événement \\ (ex: RH, Planning)};
            \node[block, right=of source] (orchestrator) {Orchestrateur \\ (Event Processor)};
            \node[db, below=of orchestrator] (db) {PostgreSQL \\ (Save-First)};
            \node[queue, right=of orchestrator] (rabbitmq) {File RabbitMQ \\ (Queue)};
            \node[block, right=of rabbitmq] (worker) {Worker \\ (Connector)};
            \node[block, below=of worker] (target) {Cible \\ (LDAP, Jira)};
            
            % Connectors
            \draw[arrow] (source) -- node[above, font=\scriptsize] {Webhook} (orchestrator);
            \draw[arrow] (orchestrator) -- node[left, font=\scriptsize] {1. Persiste} (db);
            \draw[arrow] (orchestrator) -- node[above, font=\scriptsize] {2. Publie} (rabbitmq);
            \draw[arrow] (rabbitmq) -- node[above, font=\scriptsize] {3. Dépile} (worker);
            \draw[arrow] (worker) -- node[left, font=\scriptsize] {4. Provisionne} (target);
        \end{tikzpicture}
        \caption{Flux asynchrone automatique découplé (RabbitMQ)}
        \label{fig:flow_async_hub}
    \end{figure}

    \item \textbf{Workflow synchrone manuel (mode d'administration)} : Lorsqu'un administrateur système doit intervenir en urgence depuis la console d'administration Angular (par exemple pour forcer la synchronisation d'un nouvel arrivant avec « Process Now » ou rejouer un Job en échec avec « Retry »), le Provisioning Hub bascule sur un parcours synchrone. Le \texttt{JobController} REST relaie immédiatement l'appel au \texttt{JobProcessor} qui exécute instantanément les étapes du job. La réponse REST retourne le statut d'exécution direct (succès ou échec détaillé) à la console d'administration, offrant un feedback visuel immédiat indispensable pour le support technique et la supervision en temps réel (Figure~\ref{fig:flow_sync_hub}).
    
    Plus concrètement, au sein de l'interface utilisateur (\texttt{admin-ui-angular}), l'administrateur dispose d'une console d'orchestration « Jobs \& Steps ». Si une étape échoue (statut \texttt{FAILED}), l'administrateur peut ouvrir une fenêtre modale interactive affichant le contenu exact de la requête envoyée (\texttt{reqPayload}), la réponse brute renvoyée par le connecteur cible (\texttt{resPayload}) ainsi que la trace de l'erreur (\texttt{errorMessage}). Après résolution du problème sous-jacent, l'administrateur peut relancer manuellement l'exécution de manière synchrone via l'interface, qui appelle l'API REST (\texttt{/api/v1/jobs/\{id\}/retry}) et rafraîchit immédiatement l'affichage pour confirmer la bonne exécution.
    
    \begin{figure}[H]
        \centering
        \begin{tikzpicture}[
            node distance=1.2cm and 1.8cm,
            block/.style={draw, fill=blue!10, rectangle, minimum height=0.8cm, minimum width=2.4cm, align=center, rounded corners, drop shadow, font=\small},
            arrow/.style={->, >=Stealth, thick, line width=1.2pt}
        ]
            % Nodes
            \node[block, fill=orange!15] (admin) {Administrateur \\ (UI Angular)};
            \node[block, right=of admin] (controller) {JobController \\ (API REST)};
            \node[block, right=of controller] (processor) {JobProcessor \\ (Logic Service)};
            \node[block, right=of processor, fill=red!10] (target) {LDAP / Jira \\ (Cibles Techniques)};
            
            % Connectors
            \draw[arrow] (admin.north east) -- node[above, font=\scriptsize] {1. POST} (controller.north west);
            \draw[arrow] (controller.north east) -- node[above, font=\scriptsize] {2. Appelle} (processor.north west);
            \draw[arrow] (processor.north east) -- node[above, font=\scriptsize] {3. Exécute} (target.north west);
            
            \draw[arrow, dashed] (target.south west) -- node[below, font=\scriptsize] {4. Statut} (processor.south east);
            \draw[arrow, dashed] (processor.south west) -- node[below, font=\scriptsize] {5. Réponse} (controller.south east);
            \draw[arrow, dashed] (controller.south west) -- node[below, font=\scriptsize] {6. Feedback} (admin.south east);
        \end{tikzpicture}
        \caption{Flux synchrone manuel d'administration (Feedback direct)}
        \label{fig:flow_sync_hub}
    \end{figure}
\end{itemize}

\subsection{Vue Hub 360 (Gouvernance et Audit)}
En centralisant tous les échanges de données et l'état des comptes cibles associés au sein d'un même schéma relationnel, le Provisioning Hub brise les silos de l'infrastructure héritée. L'interface d'administration offre une vue « Hub 360 » unifiée pour chaque collaborateur, permettant aux équipes de sécurité et d'audit de visualiser instantanément l'état de l'ensemble des comptes de ce collaborateur d'un seul coup d'œil.

\subsection{Moteur de Règles Déclaratif}
Grâce à l'intégration du moteur de règles SpEL (Spring Expression Language), les règles de filtrage et de routage des données sont stockées au format JSON. Un administrateur peut modifier une règle d'habilitation au runtime depuis l'interface Angular du Hub. La règle est appliquée instantanément sans nécessiter de recompilation ni de redémarrage de service.

Le Tableau~\ref{tab:comparaison_legacy_hub} résume la comparaison entre le système historique et le Provisioning Hub.

\begin{table}[H]
    \centering
    \caption{Comparaison fonctionnelle et technique : Legacy vs Provisioning Hub}
    \label{tab:comparaison_legacy_hub}
    \begin{tabular}{p{3.8cm}p{5.5cm}p{5.5cm}}
        \toprule
        \textbf{Critère} & \textbf{Système Legacy} & \textbf{Provisioning Hub} \\
        \midrule
        Mécanisme d'exécution & Scripts Batch (\texttt{.bat}) locaux et appels SOAP & Moteur événementiel (RabbitMQ) \& architecture hexagonale \\
        Mode d'exécution & Synchrone rigide ou traitements par lots isolés & Hybride : Asynchrone événementiel automatique (RabbitMQ) \& Synchrone direct pour les actions manuelles administrateur \\
        Configuration des règles & Codée en dur en Java & Déclarative (UI Angular / JSON) \\
        Prise en compte des règles & Recompilation, déploiement et redémarrage & Application instantanée en cours d'exécution via SpEL \\
        Gestion des erreurs & Arrêt silencieux sans traçabilité des scripts & Sauvegarde préalable (Save-First) et gestion de reprise (DLQ) \\
        Révocation (Offboarding) & Incomplète (génération fréquente de comptes orphelins) & Automatisée à 100\,\% dès la date de sortie RH \\
        Vue d'audit & Cloisonnée en silos de données par application & Consolidation unifiée en base PostgreSQL (Vue Hub 360) \\
        Délais de déploiement cible & Environ 15 jours de développement & Moins de 24 heures par configuration \\
        \bottomrule
    \end{tabular}
\end{table}

---

\section{Conclusion}

L'analyse de l'existant a mis en évidence le danger de s'appuyer sur des briques hétérogènes comme des scripts Batch (\texttt{.bat}) et des services SOAP synchrones pour gérer les accès de plus de 40~000 collaborateurs. Les risques de sécurité représentés par les comptes orphelins et le manque d'auditabilité dû au cloisonnement des données ont rendu indispensable la refonte complète du système. 

Le Provisioning Hub apporte une solution moderne orientée événements, résiliente et hautement configurable, garantissant une automatisation du cycle JML et une traçabilité unifiée. Le chapitre suivant détaillera les choix de conception et la modélisation architecturale adoptés pour réaliser cette vision.
 % Analyse du Contexte et Problématique
\chapter{Conception et Architecture du Système}
\label{chap:conception}

Ce chapitre présente la stratégie de conception globale et l'architecture détaillée adoptées pour le développement du Provisioning Hub (ProvisionHub). Nous y exposerons les choix structurants visant à assurer le découplage, la résilience, la sécurité et la performance du système, à travers une analyse approfondie des composants backend (Spring Boot), frontend (Angular) et de l'infrastructure de persistance et de messagerie.

\section{Introduction et Architecture Générale}
Le Provisioning Hub est conçu comme un middleware d'orchestration pour la gestion des identités et des accès (IAM). Son rôle principal est d'automatiser et de sécuriser le cycle de vie applicatif de plus de 40~000 collaborateurs au sein de l'organisation. Pour ce faire, il s'interface avec un ensemble hétérogène de référentiels sources et de systèmes cibles techniques.

L'architecture globale repose sur un paradigme hybride combinant :
\begin{itemize}
    \item \textbf{Une architecture en couches traditionnelle :} Assurant une séparation claire des responsabilités entre la présentation, l'exposition des APIs, la logique métier et la persistance.
    \item \textbf{Un découplage orienté événements (Event-Driven) :} S'appuyant sur le broker RabbitMQ pour orchestrer les traitements asynchrones et résilients.
    \item \textbf{Une architecture de connecteurs extensible (Plugin Architecture) :} Permettant d'interchanger et d'ajouter dynamiquement des modules de livraison cibles.
    \item \textbf{Un modèle CQRS (Command Query Responsibility Segregation) simplifié :} La gestion du planning des collaborateurs est déléguée en lecture temps réel au microservice distant \texttt{ms-planning-reader} (port 8082) via REST, ce qui évite la duplication locale d'une base de données de plannings volumineuse, tout en maintenant les écritures et l'agrégation dans le cœur du Provisioning Hub pour la vue consolidée à 360 degrés.
\end{itemize}

Le schéma d'architecture globale est représenté dans la Figure~\ref{fig:architecture_globale}.

\begin{figure}[H]
    \centering
    \begin{tikzpicture}[
        node distance=1.2cm and 1.8cm,
        block/.style={draw, fill=blue!10, rectangle, minimum height=0.8cm, minimum width=2.8cm, align=center, rounded corners, drop shadow, font=\small},
        database/.style={draw, cylinder, fill=gray!20, minimum height=1.0cm, minimum width=1.6cm, shape border rotate=90, align=center, drop shadow, font=\small},
        arrow/.style={->, >=Stealth, thick, line width=1.2pt}
    ]
        % Nodes
        \node[block, fill=blue!15] (angular) {Interface Admin\\(Angular SPA)};
        \node[block, below=of angular, fill=orange!15] (backend) {Backend API Core\\(Spring Boot)};
        \node[database, below=of backend, fill=yellow!15] (postgres) {Base PostgreSQL\\(PostgreSQL 17)};
        \node[block, right=of backend, fill=red!15] (rabbitmq) {Message Bus\\(RabbitMQ)};
        \node[block, right=of rabbitmq, fill=green!15] (connectors) {Connecteurs cibles\\(LDAP, Jira, Gmail)};
        \node[block, right=of angular, fill=purple!15] (planning) {Planning Service\\(ms-planning-reader)};

        % Connectors
        \draw[arrow] (angular) -- node[left, font=\scriptsize] {REST API} (backend);
        \draw[arrow] (backend) -- node[left, font=\scriptsize] {JPA / SQL} (postgres);
        \draw[arrow] (backend) -- node[above, font=\scriptsize] {AMQP Publish} (rabbitmq);
        \draw[arrow] (rabbitmq) -- node[above, font=\scriptsize] {Events Dispatch} (connectors);
        \draw[arrow] (angular) -- node[above, font=\scriptsize] {REST (Port 8082)} (planning);
        \draw[arrow, dashed] (backend) -- node[right, font=\scriptsize] {CQRS REST} (planning);
    \end{tikzpicture}
    \caption{Architecture globale en couches et flux de données du Provisioning Hub}
    \label{fig:architecture_globale}
\end{figure}

\section{Stack Technologique}
Le choix de la stack technologique répond à des contraintes industrielles de maintenabilité, d'évolutivité et de support communautaire. Les tableaux suivants résument les composants logiciels exploités.

\subsection{Composants Backend}
Le backend s'appuie sur l'écosystème Java/Spring, comme détaillé dans le Tableau~\ref{tab:stack_backend}.

\begin{table}[H]
    \centering
    \caption{Stack technologique et versions logicielles du Backend}
    \label{tab:stack_backend}
    \begin{tabular}{lll}
        \toprule
        \textbf{Technologie} & \textbf{Version} & \textbf{Rôle et justification} \\
        \midrule
        Java & 21 (LTS) & Support des fonctionnalités modernes et des threads virtuels. \\
        Spring Boot & 3.3.5 & Framework d'infrastructures et de démarrage rapide. \\
        Spring Data JPA & Inclus & Abstraction de la couche d'accès aux données (ORM). \\
        Spring Security & Inclus & Sécurisation stateless et gestion des politiques RBAC. \\
        Spring AMQP & Inclus & Intégration native avec le broker RabbitMQ. \\
        Hibernate & 6.3 & Implémentation du standard de persistance JPA. \\
        Hypersistence Utils & 3.8.3 & Prise en charge performante du type JSONB de PostgreSQL. \\
        Flyway & Inclus & Versioning et migration automatique du schéma relationnel. \\
        Lombok & Inclus & Réduction du code boilerplate (générateurs automatiques). \\
        \bottomrule
    \end{tabular}
\end{table}

\subsection{Composants Frontend}
La console d'administration est développée sous forme d'application monopage (SPA) avec le framework Angular (Tableau~\ref{tab:stack_frontend}).

\begin{table}[H]
    \centering
    \caption{Stack technologique et versions logicielles du Frontend}
    \label{tab:stack_frontend}
    \begin{tabular}{lll}
        \toprule
        \textbf{Technologie} & \textbf{Version} & \textbf{Rôle et justification} \\
        \midrule
        Angular & 17.3.0 & Framework de composants moderne (Standalone Components). \\
        TypeScript & 5.4.0 & Typage statique fort pour fiabiliser la logique applicative. \\
        RxJS & 7.8.0 & Gestion réactive des flux d'événements et requêtes HTTP. \\
        SCSS & Natif & Préprocesseur CSS pour structurer le design system. \\
        \bottomrule
    \end{tabular}
\end{table}

\subsection{Infrastructure et Déploiement}
L'infrastructure cible repose sur des conteneurs isolés (Tableau~\ref{tab:stack_infra}).

\begin{table}[H]
    \centering
    \caption{Composants d'infrastructure et d'orchestration}
    \label{tab:stack_infra}
    \begin{tabular}{lll}
        \toprule
        \textbf{Composant} & \textbf{Version} & \textbf{Rôle opérationnel} \\
        \midrule
        PostgreSQL & 17-alpine & Base de données relationnelle et transactionnelle. \\
        RabbitMQ & 3.13-management & Gestionnaire de files d'attente et bus de messages. \\
        Docker Compose & 3.8 & Description et orchestration locale des conteneurs. \\
        Nginx & Alpine & Serveur web pour le frontend et reverse-proxy de routage. \\
        \bottomrule
    \end{tabular}
\end{table}

\section{Conception Modulaire Backend (Spring Boot)}
Le code source du backend est structuré selon une arborescence modulaire stricte afin d'isoler les responsabilités et de faciliter la maintenance évolutive.

\subsection{Structure des Packages}
Les packages principaux du projet s'organisent de la manière suivante :
\begin{itemize}
    \item \texttt{com.company.middleware.config} : Configuration système (JPA, DataSources multiples pour PostgreSQL et le portail HR MySQL historique, sécurité, RabbitMQ, schedulers).
    \item \texttt{com.company.middleware.entity} : Modélisation des 37 entités JPA du système.
    \item \texttt{com.company.middleware.repository} : Interfaces Spring Data JPA pour l'accès aux données.
    \item \texttt{com.company.middleware.service} : Implémentations des règles métiers (gestion des employés, logs, alertes).
    \item \texttt{com.company.middleware.controller} : Contrôleurs REST exposant les points d'entrée de l'API.
    \item \texttt{com.company.middleware.orchestration} : Moteur d'évaluation des règles, ordonnanceur de jobs et gestionnaire de priorités.
    \item \texttt{com.company.middleware.connector} : Définition de l'interface commune des connecteurs et implémentations physiques (LDAP, Jira, Mail).
    \item \texttt{com.company.middleware.messaging} : Composants d'émission et de réception RabbitMQ.
\end{itemize}

\subsection{Statistiques de Code}
Afin d'illustrer la complexité du socle applicatif du middleware, le Tableau~\ref{tab:statistiques_code} présente la répartition des classes clés.

\begin{table}[H]
    \centering
    \caption{Volume et typologie des classes du backend Spring Boot}
    \label{tab:statistiques_code}
    \begin{tabular}{lc}
        \toprule
        \textbf{Type de composant Java} & \textbf{Nombre approximatif de classes} \\
        \midrule
        Entités JPA & 37 \\
        Repositories Spring Data & 37 \\
        Services métiers & 18 \\
        Contrôleurs REST (dont webhooks d'intégration) & 25 \\
        DTOs (Data Transfer Objects) & 16 \\
        Composants de l'orchestrateur & 11 \\
        Composants du moteur de workflows et d'échange & 11 \\
        Connecteurs et services techniques associés & 6 \\
        \bottomrule
    \end{tabular}
\end{table}

\section{Conception du Client Frontend (Angular)}
L'interface utilisateur du Provisioning Hub sert de console d'administration pour la configuration opérationnelle et le suivi en temps réel des transactions.

\subsection{Modèle Standalone et Chargement Différé}
Le projet Angular s'aligne sur les standards d'Angular 17 en abandonnant l'usage des modules (\texttt{NgModules}) au profit des \textit{Standalone Components}. Chaque vue ou sous-composant déclare directement ses dépendances, ce qui simplifie le graphe de compilation. 
De plus, le routage de l'application utilise systématiquement le chargement différé (\textit{Lazy Loading}) via des promesses dynamiques d'importation. Cela garantit que seul le code requis pour la page consultée est transféré au navigateur, optimisant ainsi drastiquement les performances d'affichage et de réactivité.

\subsection{Structure et Architecture de Navigation}
L'application se découpe en trois dossiers structurants sous \texttt{src/app/} :
\begin{itemize}
    \item \texttt{core/} : Contient les services d'infrastructure transversaux, tels que \texttt{AuthService} (gestion des sessions RBAC), \texttt{ToastService} (notifications), et \texttt{ApiService}.
    \item \texttt{shared/} : Regroupe les composants réutilisables à travers l'application (modales, badges de statuts, barres de pagination).
    \item \texttt{pages/} : Implémente l'ensemble des 21 écrans fonctionnels (dashboard, profils des employés, édition de politiques IAM, monitoring de la DLQ, logs d'échanges).
\end{itemize}

\subsection{Centralisation des Échanges HTTP}
Pour interagir avec le backend, le frontend s'appuie sur le patron de conception Facade à travers la classe \texttt{ApiService}. Ce service encapsule l'ensemble des requêtes REST (plus de 80 méthodes de requêtes typées) en masquant la manipulation des en-têtes d'authentification et les transformations RxJS sous-jacentes.

\section{Modélisation et Persistance des Données (PostgreSQL 17)}
La modélisation des données a été pensée pour répondre à deux exigences majeures : la traçabilité absolue de chaque décision d'accès et le support d'importants volumes d'échanges quotidiens.

\subsection{Entités Principales et Relations}
Le schéma relationnel est centré autour de l'entité \texttt{Employee}. Un employé possède un état cible (\texttt{EmployeeTargetState}) calculé pour chaque connecteur, un historique de modifications (\texttt{EmployeeHistory}), et peut posséder des enregistrements d'événements à 360 degrés (congés, plannings, logs SSO). 
Par ailleurs, l'exécution d'un changement déclenche un \texttt{Job} composé de plusieurs \texttt{JobStep} (étapes unitaires idempotentes) pouvant générer des erreurs collectées dans la table \texttt{DeadLetter}. Enfin, les politiques de provisioning (\texttt{Policy}) sont rattachées aux connecteurs.

\subsection{Optimisations Avancées de PostgreSQL 17}
Pour maintenir des temps de réponse faibles malgré l'accumulation des logs de synchronisation, plusieurs techniques avancées sont configurées en base de données :
\begin{itemize}
    \item \textbf{Partitionnement horizontal (Table Partitioning) :} Afin d'éviter la dégradation des index sur les tables à forte croissance, six tables clés sont partitionnées par plage de dates (\texttt{PARTITION BY RANGE}) sur les années civiles : \texttt{employee\_history}, \texttt{job}, \texttt{step\_log}, \texttt{decision\_audit}, \texttt{state\_history} et \texttt{admin\_audit}.
    \item \textbf{Champs extensibles JSONB et index GIN :} Pour permettre aux administrateurs d'intégrer de nouveaux champs collaborateurs sans altérer le schéma SQL, l'entité \texttt{Employee} dispose d'une colonne JSONB nommée \texttt{attrs}. L'accès à ces attributs dynamiques est optimisé par la création d'index GIN (\textit{Generalized Inverted Index}) s'appuyant sur l'extension \texttt{pg\_trgm}.
    \item \textbf{Garantie d'idempotence :} Pour éviter les exécutions multiples d'un même ordre de provisionnement en cas de perturbation réseau, un index d'unicité est positionné sur la colonne \texttt{idempotency\_key} de la table \texttt{job\_step}.
    \item \textbf{Vues de supervision opérationnelle :} Plusieurs requêtes d'agrégation complexes sont encapsulées dans des vues SQL natives (\texttt{v\_backlog}, \texttt{v\_dlq\_open}, \texttt{v\_drift\_open}) afin de simplifier les requêtes du dashboard.
\end{itemize}

\section{Conception des APIs REST}
L'exposition des fonctionnalités du Provisioning Hub s'effectue au moyen d'APIs REST normalisées.

\subsection{Conventions de Conception}
\begin{itemize}
    \item \textbf{Versioning de l'API :} Les routes métiers de base sont préfixées par \texttt{/api/v1/}, tandis que les routes associées au nouveau moteur d'ingestion dynamique (Extract-Load-Transform) sont exposées sous \texttt{/api/v2/sources} et \texttt{/api/v2/destinations}.
    \item \textbf{Stateless et Pagination :} Les requêtes de lecture retournent des enveloppes paginées standardisées basées sur le modèle \texttt{PageResponse<T>} de Spring Data, évitant ainsi le transfert inutile de collections massives.
    \item \textbf{Traitement synchrone et asynchrone :} Les actions de mise à jour de configuration s'exécutent de façon synchrone. En revanche, le déclenchement d'un provisionnement renvoie immédiatement un statut \texttt{202 Accepted} contenant l'identifiant du job, dont l'exécution se poursuit de manière asynchrone en tâche de fond.
\end{itemize}

\subsection{Catalogue des Endpoints Clés}
Le Tableau~\ref{tab:catalogue_endpoints} liste les points d'accès critiques de l'API REST du Provisioning Hub.

\begin{table}[H]
    \centering
    \caption{Endpoints clés de l'API REST}
    \label{tab:catalogue_endpoints}
    \begin{tabular}{lcl}
        \toprule
        \textbf{Route} & \textbf{Méthode} & \textbf{Rôle fonctionnel} \\
        \midrule
        \texttt{/api/v1/employees} & GET & Recherche paginée de collaborateurs. \\
        \texttt{/api/v1/employees/\{id\}/360} & GET & Consolidation de la vue 360 degrés. \\
        \texttt{/api/v1/connectors/\{key\}/health} & GET & Diagnostic de santé d'un système cible. \\
        \texttt{/api/v1/policies} & POST & Enregistrement d'une nouvelle règle IAM. \\
        \texttt{/api/v1/jobs/backfill} & POST & Lancement d'une synchronisation de masse. \\
        \texttt{/api/v1/dlq/\{id\}/retry} & POST & Rejeu d'un message en erreur depuis la DLQ. \\
        \texttt{/api/v2/sources/\{key\}/trigger} & POST & Lancement manuel de l'ingestion ETL. \\
        \bottomrule
    \end{tabular}
\end{table}

\section{Moteur de Politiques IAM (Policy Engine)}
Le moteur de politiques est le cœur logique qui détermine le droit d'accès et le formatage des comptes des collaborateurs sur chaque système cible.

\subsection{Évaluation Dynamique avec SpEL}
Pour éviter de coder les règles métier en dur dans l'application, le Provisioning Hub intègre le moteur Spring Expression Language (SpEL). Les expressions sont stockées sous forme de chaînes de caractères en base de données et évaluées au runtime contre l'objet \texttt{Employee} en cours de traitement. 
Ainsi, une expression telle que \texttt{employee.status == 'ACTIVE' \&\& employee.country == 'FR'} sera évaluée à \texttt{true} ou \texttt{false} de manière dynamique, modifiant instantanément le droit d'accès de l'utilisateur sans requérir de redémarrage de l'application.

\subsection{Règles Logiques JSON et Évaluateurs Spécialisés}
En complément de SpEL, le moteur gère des arbres de conditions logiques au format JSON, supportant les opérateurs logiques complexes (groupements \texttt{all} pour AND, et \texttt{any} pour OR) associés à des comparateurs typés (\texttt{EQUALS}, \texttt{CONTAINS}, \texttt{REGEX}, \texttt{BEFORE}, etc.). 
Pour les cas d'évaluation trop complexes ou nécessitant des requêtes réseau, l'interface \texttt{CustomPolicyEvaluator} permet de développer des classes spécifiques (comme \texttt{LdapCustomPolicyEvaluator}) qui sont résolues dynamiquement grâce à un registre d'évaluateurs (\texttt{CustomPolicyEvaluatorRegistry}).

L'algorithme de décision d'évaluation des règles est représenté dans la Figure~\ref{fig:algorithme_policy}.

\begin{figure}[H]
    \centering
    \begin{tikzpicture}[
        node distance=1.2cm and 1.5cm,
        decision/.style={draw, fill=yellow!10, diamond, aspect=1.8, minimum width=2.5cm, align=center, font=\scriptsize},
        block/.style={draw, fill=blue!10, rectangle, minimum height=0.8cm, minimum width=2.8cm, align=center, rounded corners, drop shadow, font=\small},
        arrow/.style={->, >=Stealth, thick, line width=1.2pt}
    ]
        % Nodes
        \node[block] (start) {Évaluation Collaborateur\\(Connector)};
        \node[block, below=of start] (load) {Charger règles actives\\par priorité};
        \node[decision, below=of load] (check_custom) {Évaluateur\\Spécialisé ?};
        \node[block, right=of check_custom, xshift=1cm] (exec_custom) {Exécuter\\Java Custom Evaluator};
        \node[block, below=of check_custom] (exec_spel) {Évaluer expressions\\SpEL / Conditions JSON};
        \node[decision, below=of exec_spel] (match) {Condition\\validée ?};
        \node[block, below=of match] (apply) {Appliquer Action\\(ALLOW / BLOCK / SKIP)};
        \node[block, left=of match, xshift=-1cm] (next) {Passer à la\\règle suivante};

        % Connectors
        \draw[arrow] (start) -- (load);
        \draw[arrow] (load) -- (check_custom);
        \draw[arrow] (check_custom) -- node[above, font=\scriptsize] {Oui} (exec_custom);
        \draw[arrow] (check_custom) -- node[right, font=\scriptsize] {Non} (exec_spel);
        \draw[arrow] (exec_custom.south) |- (match.east);
        \draw[arrow] (exec_spel) -- (match);
        \draw[arrow] (match) -- node[right, font=\scriptsize] {Oui} (apply);
        \draw[arrow] (match) -- node[above, font=\scriptsize] {Non} (next);
        \draw[arrow] (next) |- (load.west);
    \end{tikzpicture}
    \caption{Algorithme de décision du moteur d'évaluation des politiques IAM}
    \label{fig:algorithme_policy}
\end{figure}

\section{Architecture Event-Driven (RabbitMQ)}
L'orchestration asynchrone du provisionnement repose sur un réseau de files d'attente configuré au sein du broker de messages RabbitMQ.

\subsection{Topologie de Messagerie}
Le système configure un Topic Exchange principal nommé \texttt{provisionhub.exchange}. Les événements relatifs aux collaborateurs (créations, mutations, départs) y sont publiés avec des clés de routage structurées, par exemple \texttt{employee.event.created} ou \texttt{employee.assignment.changed}. 
Grâce à ce mécanisme, plusieurs files d'attente s'abonnent à des sous-ensembles d'événements à l'aide de wildcards (\texttt{employee.\#}), assurant un routage sélectif et parallèle des messages vers les différents sous-systèmes d'intégration.

\subsection{Haute Fiabilité et Mécanisme de DLQ}
Afin de prévenir toute perte d'événements en cas d'indisponibilité d'un système cible ou du backend, les règles suivantes sont appliquées au niveau de RabbitMQ :
\begin{itemize}
    \item \textbf{Durabilité :} L'exchange et l'ensemble des files d'attente sont configurés comme \texttt{durable}, ce qui signifie qu'ils sont persistés sur disque et résistent à un redémarrage du broker.
    \item \textbf{Publisher Confirms :} Le backend attend un accusé de réception de RabbitMQ avant de considérer un événement comme transmis.
    \item \textbf{Dead Letter Queues (DLQ) :} Chaque file d'attente principale est associée à une file de rejet (\texttt{Dead Letter Queue}). En cas d'erreur de traitement persistante sur un message (après épuisement des tentatives de rejeu avec délai exponentiel), RabbitMQ déplace automatiquement le message dans la DLQ correspondante. Ces anomalies sont historisées dans la table \texttt{DeadLetter} et exposées dans l'UI pour rejeu ou suppression par un administrateur.
\end{itemize}

Le flux de routage des événements et de traitement des rejets est schématisé dans la Figure~\ref{fig:flux_rabbitmq}.

\begin{figure}[H]
    \centering
    \begin{tikzpicture}[
        node distance=1.1cm and 1.6cm,
        block/.style={draw, fill=blue!10, rectangle, minimum height=0.8cm, minimum width=2.8cm, align=center, rounded corners, drop shadow, font=\small},
        queue/.style={draw, fill=yellow!15, rectangle, minimum height=0.8cm, minimum width=2.6cm, align=center, drop shadow, font=\small},
        arrow/.style={->, >=Stealth, thick, line width=1.2pt}
    ]
        % Nodes
        \node[block] (pub) {Producteur d'événements\\(HubEventPublisher)};
        \node[block, right=of pub, fill=orange!15] (exch) {Topic Exchange\\(provisionhub.exchange)};
        \node[queue, below=of exch] (q_main) {File Principale\\(employee-events)};
        \node[block, left=of q_main] (listener) {Consommateur\\(RabbitMQEventListener)};
        \node[queue, below=of q_main, fill=red!10] (q_dlq) {Dead Letter Queue\\(employee-events.dlq)};
        \node[block, left=of q_dlq, fill=green!15] (dlq_service) {Console DLQ\\(Rejeu manuel / Diagnostic)};

        % Connectors
        \draw[arrow] (pub) -- node[above, font=\scriptsize] {Publish} (exch);
        \draw[arrow] (exch) -- node[right, font=\scriptsize] {routing: employee.\#} (q_main);
        \draw[arrow] (q_main) -- node[above, font=\scriptsize] {Consume} (listener);
        \draw[arrow] (listener) -- node[left, font=\scriptsize] {Erreur persistante} (q_dlq);
        \draw[arrow] (q_dlq) -- node[above, font=\scriptsize] {Inspecter} (dlq_service);
        \draw[arrow, dashed] (dlq_service) -- node[left, font=\scriptsize] {Rejeu} (exch);
    \end{tikzpicture}
    \caption{Topologie de routage des messages et cycle de vie des erreurs dans RabbitMQ}
    \label{fig:flux_rabbitmq}
\end{figure}

\section{Système de Connecteurs (Plugin Architecture)}
Le provisionnement effectif vers les systèmes d'information cibles s'effectue via des modules spécialisés appelés connecteurs.

\subsection{Contrat d'Interface et Strategy Pattern}
Pour éviter le couplage fort entre le moteur de règles et la tuyauterie technique propre à chaque destination, nous avons implémenté le patron de conception Strategy. Chaque type d'intégration implémente l'interface commune \texttt{ConnectorPlugin}. Celle-ci définit un ensemble d'opérations normalisées :
\begin{itemize}
    \item \texttt{testConnection()} : Vérification de la disponibilité du système distant.
    \item \texttt{createUser()} / \texttt{updateUser()} : Manipulation des fiches utilisateurs.
    \item \texttt{disableUser()} / \texttt{deleteUser()} : Actions de désactivation temporaire ou de déprovisionnement.
    \item \texttt{readUser()} / \texttt{getUserGroups()} : Lecture de l'état distant pour la détection de dérive (\textit{Drift Detection}).
\end{itemize}

\subsection{Auto-Découverte via le ConnectorRegistry}
L'ajout d'un nouveau connecteur s'effectue par simple injection de sa classe implémentant \texttt{ConnectorPlugin} dans le conteneur Spring. Au démarrage de l'application, la classe \texttt{ConnectorRegistry} scanne automatiquement le contexte d'application grâce à la méthode d'initialisation \texttt{@PostConstruct}. Les plug-ins détectés sont référencés dans une map interne indexée par leur type de protocole (ex : \texttt{LDAP}, \texttt{REST\_API}). Le système peut ainsi instancier et appeler dynamiquement le bon connecteur lors du traitement d'une étape de job.

\section{Sécurité et Contrôle d'Accès (RBAC)}
En manipulant des données d'identités critiques et des secrets de connexion, le Provisioning Hub applique des règles strictes de sécurité.

\subsection{Sécurité Backend}
Le framework Spring Security gère la protection du backend :
\begin{itemize}
    \item \textbf{Authentification Stateless :} Aucun état de session n'est maintenu sur le serveur. Chaque requête HTTP doit porter un jeton d'authentification valide.
    \item \textbf{Protection CSRF et CORS :} La protection CSRF est désactivée car l'API est sans état. Les règles CORS (\textit{Cross-Origin Resource Sharing}) sont configurées pour n'autoriser les appels que depuis le domaine légitime hébergeant la console d'administration Angular.
\end{itemize}

\subsection{Modèle RBAC Frontend}
Le contrôle d'accès basé sur les rôles (RBAC) est géré côté frontend pour restreindre l'affichage et l'interaction avec les fonctionnalités selon le profil de l'utilisateur connecté. Quatre rôles hiérarchiques sont configurés :
\begin{enumerate}
    \item \textbf{SUPER\_ADMIN :} Accès complet à l'ensemble du système, y compris la configuration de la sécurité et la rotation des secrets.
    \item \textbf{OPERATOR :} Rôle de supervision opérationnelle. Il possède les droits d'administration des employés, de rejeu de la DLQ, de gestion des connecteurs et de configuration des sources.
    \item \textbf{AUDITOR :} Rôle d'inspection. Il peut visualiser le dashboard, la timeline des employés, les alertes de dérives et générer des rapports d'audit.
    \item \textbf{VIEWER :} Profil en lecture seule. Il peut uniquement consulter le dashboard et la fiche de synthèse des collaborateurs.
\end{enumerate}

Le flux de navigation du frontend est protégé par des gardes d'activation de routes (\texttt{AuthGuard}) qui interceptent les transitions et vérifient les permissions associées au rôle stocké de manière sécurisée en session.

\section{Déploiement et Architecture Conteneurisée (Docker)}
Pour simplifier l'exploitation et assurer l'isolation des différents modules techniques, l'application est packagée sous forme de microservices conteneurisés configurés via Docker Compose.

Le réseau virtuel \texttt{provisionhub-network} isole la stack logicielle. Les appels externes sont filtrés et routés par Nginx. Les conteneurs du cœur applicatif (\texttt{ProvisionHub} et \texttt{ms-hr-collector}) communiquent en interne avec PostgreSQL et RabbitMQ sans exposer leurs ports sur le réseau public. De même, les microservices des connecteurs cibles s'exécutent dans des conteneurs dédiés, ce qui permet de faire évoluer ou de redémarrer un connecteur spécifique (ex : \texttt{ms-connector-jira}) sans impacter le fonctionnement général de la plateforme.

La topologie du réseau et du déploiement Docker Compose est illustrée dans la Figure~\ref{fig:docker_network}.

\begin{figure}[H]
    \centering
    \begin{tikzpicture}[
        node distance=1.1cm and 1.3cm,
        container/.style={draw, fill=blue!10, rectangle, minimum height=0.7cm, minimum width=2.4cm, align=center, rounded corners, drop shadow, font=\scriptsize},
        boundary/.style={draw, dashed, inner sep=0.4cm, label={[anchor=south, font=\small]north:Réseau Bridge : provisionhub-network}},
        arrow/.style={->, >=Stealth, thick, line width=1.0pt}
    ]
        % Docker network boundaries
        \begin{scope}
            \node[container, fill=orange!15] (core) {ProvisionHub\\(Core Engine)};
            \node[container, below=of core, fill=yellow!15] (db) {PostgreSQL 17};
            \node[container, right=of core, fill=red!15] (mq) {RabbitMQ};
            \node[container, left=of core, fill=purple!15] (coll) {ms-hr-collector};
            \node[container, below=of mq, fill=green!15] (conn1) {ms-connector-ldap};
            \node[container, right=of mq, fill=green!15] (conn2) {ms-connector-jira};
        \end{scope}

        % Grouping into a dotted box representing the docker network
        \node[boundary, fit=(core) (db) (mq) (coll) (conn1) (conn2)] (network) {};

        % External nodes outside the docker network
        \node[container, above=of core, fill=blue!15, yshift=0.5cm] (nginx) {Reverse Proxy\\(Nginx / Static UI)};

        % Connections
        \draw[arrow] (nginx) -- node[right, font=\tiny] {Port 8080} (core);
        \draw[arrow] (core) -- (db);
        \draw[arrow] (core) -- (mq);
        \draw[arrow] (coll) -- (mq);
        \draw[arrow] (mq) -- (conn1);
        \draw[arrow] (mq) -- (conn2);
    \end{tikzpicture}
    \caption{Réseau virtuel Docker Compose et isolation des microservices applicatifs}
    \label{fig:docker_network}
\end{figure}

\section{Patrons de Conception (Design Patterns) Appliqués}
Le Provisioning Hub fait un usage intensif de patrons de conception éprouvés pour résoudre des problématiques d'architecture logicielle récurrentes.

\subsection{Patrons Architecturaux}
\begin{itemize}
    \item \textbf{Layered Architecture (Architecture en couches) :} Découpage horizontal en Contrôleurs $\rightarrow$ Services $\rightarrow$ Repositories.
    \item \textbf{Event-Driven Architecture (EDA) :} Communication par messages asynchrones via RabbitMQ pour découpler la détection du changement et le provisionnement.
    \item \textbf{CQRS (Command Query Responsibility Segregation) simplifié :} Séparation logique des flux de mise à jour et d'extraction de données (notamment pour les plannings du module 360°).
\end{itemize}

\subsection{Patrons de Conception GoF (Gang of Four)}
\begin{itemize}
    \item \textbf{Strategy Pattern :} Implémenté pour les connecteurs cibles via l'interface commune \texttt{ConnectorPlugin}, permettant de rendre interchangeable la logique d'appel technique.
    \item \textbf{Registry Pattern :} Encapsulé dans \texttt{ConnectorRegistry} et \texttt{CustomPolicyEvaluatorRegistry} pour maintenir la liste des services disponibles découverts par réflexion.
    \item \textbf{Observer Pattern :} Matérialisé par \texttt{RabbitMQEventListener} qui écoute de façon non bloquante l'apparition d'événements sur le bus RabbitMQ.
    \item \textbf{Template Method Pattern :} Configuré dans la classe d'ordonnancement \texttt{JobProcessor.processNextBatch()} qui déroule un squelette d'exécution fixe (extraction, verrouillage pessimiste, boucle d'évaluation, journalisation) tout en déléguant la logique de mapping à des sous-composants.
    \item \textbf{Builder Pattern :} Appliqué sur l'intégralité des DTOs et entités d'échange via l'annotation \texttt{@Builder} de Lombok afin de fiabiliser la construction des payloads d'API.
\end{itemize}

\subsection{Patrons de Résilience et d'Infrastructure}
\begin{itemize}
    \item \textbf{Dead Letter Queue (DLQ) :} Routage automatique des messages en échec persistant pour éviter de bloquer la file principale et permettre un audit.
    \item \textbf{Retry with Exponential Backoff (Reprise avec délai exponentiel) :} Algorithme géré par le \texttt{RetryManager} pour espacer progressivement les tentatives d'appels vers un système tiers en panne (ex : délai de 1s, 5s, 25s...).
    \item \textbf{Pessimistic Locking (SELECT FOR UPDATE SKIP LOCKED) :} Utilisé dans la requête SQL de chargement des jobs par le \texttt{JobProcessor} pour autoriser une exécution distribuée sur plusieurs nœuds sans conflit de double traitement.
    \item \textbf{Drift Detection (Détection de dérive) :} Script autonome du \texttt{DriftScanner} qui compare à intervalle fixe l'état théorique stocké en base de données et l'état réel lu sur la cible LDAP/Jira afin de lever des alertes d'anomalies.
\end{itemize}

\section{Spécifications des Diagrammes UML à Produire}
Pour guider la compréhension visuelle de la modélisation et de la dynamique du système, le rapport de PFE s'appuiera sur la production des diagrammes UML détaillés ci-dessous :

\begin{enumerate}
    \item \textbf{Diagramme de Cas d'Utilisation (Use Case Diagram) :} Représentant les interactions des rôles (Super Admin, Opérateur, Auditeur) avec les fonctionnalités majeures (supervision, rejeu de DLQ, test de connexion, gestion des politiques).
    \item \textbf{Diagramme de Classes Global (Class Diagram) :} Modélisant la structure statique des entités métiers (\texttt{Employee}, \texttt{Job}, \texttt{JobStep}, \texttt{DeadLetter}, \texttt{Policy}, \texttt{Connector}) et leurs cardinalités.
    \item \textbf{Diagrammes de Séquence Systèmes (Sequence Diagrams) :}
        \begin{itemize}
            \item Le scénario de \textit{Provisioning} classique, décrivant les échanges chronologiques entre le webhook d'intégration, l'orchestrateur de jobs, le moteur de politiques et l'appel de connecteur.
            \item Le scénario d'\textit{Agrégation 360°}, illustrant la capture d'événements (badgeuse, congés) et leur acheminement par RabbitMQ vers le service d'agrégation.
        \end{itemize}
    \item \textbf{Diagramme d'Activité (Activity Diagram) :} Représentant graphiquement le flux logique séquentiel d'évaluation du moteur de règles.
    \item \textbf{Diagramme de Déploiement (Deployment Diagram) :} Cartographiant l'environnement de production composé de l'instance PostgreSQL, du broker RabbitMQ, et des conteneurs Spring Boot et Angular isolés sur le serveur d'intégration.
\end{enumerate}

\section{Conclusion}
La conception et l'architecture du Provisioning Hub ont été structurées pour offrir une modularité et une extensibilité optimales. En dissociant les responsabilités applicatives à travers une architecture en couches, en fiabilisant les communications réseau au moyen de files d'attente asynchrones, et en adoptant une modélisation relationnelle robuste sous PostgreSQL 17, nous avons jeté les bases d'un middleware résilient et performant. Le chapitre suivant décrira l'implémentation technique et les défis de programmation relevés au cours du développement.
 % Conception et Architecture du Système
\chapter{Réalisation Technique et Implémentation}
\label{chap:realisation}

Ce chapitre détaille la mise en œuvre technique du Provisioning Hub. Nous passerons en revue la stack technologique, l'implémentation des patterns de résilience et les optimisations de performance qui permettent au système de traiter des volumes de données importants avec une fiabilité maximale.

\section{Introduction}
L'implémentation d'une plateforme déclarative demande une attention particulière sur la gestion du cycle de vie des données et la robustesse des communications réseau. L'enjeu est de transformer des configurations abstraites en appels API concrets, sécurisés et auditables.

\section{Environnement de Développement}
\label{sec:environnement_dev}

\subsection{Stack Technique et Justifications}
Le choix des technologies a été guidé par des impératifs de performance, de maintenabilité et de support à long terme.
\begin{itemize}
    \item \textbf{Java 21 \& Spring Boot 3.3 :} L'utilisation de la version 21 permet de bénéficier des \textit{Virtual Threads} (Project Loom), réduisant drastiquement la consommation de ressources lors des appels réseau bloquants vers les destinations. Spring Boot 3.3 assure une intégration fluide des modules de sécurité et de persistance.
    \item \textbf{Angular 19 :} Choisi pour sa capacité à gérer des états complexes dans l'interface d'administration, notamment pour le concepteur de workflows et la vue Hub 360.
    \item \textbf{RabbitMQ 3.13 :} Assure le rôle de broker de messages, permettant l'implémentation d'un mode de traitement asynchrone, la gestion des priorités et la mise en place de files d'attente de récupération.
    \item \textbf{PostgreSQL 17 :} Offre des performances accrues pour les requêtes complexes et une fiabilité transactionnelle indispensable pour le pattern Save-First.
\end{itemize}

\subsection{Orchestration et CI/CD}
Pour garantir la portabilité et la reproductibilité du système, nous avons utilisé \textbf{Docker} pour la conteneurisation de l'application et de ses dépendances. La gestion des schémas de base de données est assurée par \textbf{Flyway}, permettant un versioning rigoureux des migrations SQL. Ce flux CI/CD permet d'assurer que chaque modification de schéma est testée et appliquée de manière atomique sur tous les environnements.

\section{Implémentation des Fonctionnalités Clés}
\label{sec:impl_cles}

\subsection{Le pattern "Save-First" et la traçabilité absolue}
L'un des piliers du système est la garantie qu'aucune donnée n'est perdue, même en cas de crash système. Nous avons implémenté le pattern \textbf{Save-First}. Le flux d'exécution suit strictement cet ordre :
\begin{enumerate}
    \item \textbf{Construction :} Le système assemble le payload final en fonction des règles de transformation.
    \item \textbf{Persistance :} L'échange est enregistré en base de données avec l'état \texttt{PENDING} et le contenu du payload.
    \item \textbf{Transmission :} L'appel réseau est effectué vers la destination.
    \item \textbf{Finalisation :} L'état de l'échange est mis à jour (\texttt{SUCCESS} ou \texttt{FAILED}) avec la réponse brute de la cible.
\end{enumerate}
Cette approche transforme la base de données en un journal d'audit immuable, permettant la reconstruction complète de l'état du système à tout instant.

\subsection{Gestion de la résilience et des erreurs}
La communication avec des systèmes externes est intrinsèquement instable. Pour y remédier, nous avons mis en place :
\begin{itemize}
    \item \textbf{Retry avec Backoff Exponentiel :} En cas d'erreur transitoire (ex: timeout, 503 Service Unavailable), le système reprogramme l'échange avec un délai croissant (ex: 1s, 10s, 1min, 10min). Cela évite de saturer une cible déjà en difficulté.
    \item \textbf{Dead Letter Queue (DLQ) :} Si le nombre maximum de tentatives est atteint, le message est déplacé vers une DLQ. Cette file d'attente sert de zone d'isolement. L'administrateur peut alors analyser l'erreur via l'interface, corriger la règle métier et déclencher un \textit{replay} manuel de l'échange.
\end{itemize}

\subsection{Optimisation des performances et résolution du problème N+1}
Le traitement de volumes massifs de données (plusieurs dizaines de milliers de collaborateurs) a révélé un goulot d'étranglement lié aux requêtes répétitives vers la base de données (problème N+1). Nous avons optimisé ce processus via le \textbf{Bulk Loading}. 
Au lieu de récupérer les attributs pour chaque collaborateur individuellement, le système récupère les données par lots (\textit{batching}) en utilisant des jointures optimisées et le chargement anticipé (\textit{Eager Loading}) via JPA. Cette optimisation a permis de réduire le temps de traitement global de 70\%, augmentant ainsi considérablement le débit d'ingestion.

\section{Sécurité et Conformité}
\label{sec:securite}

La manipulation de données RH et de secrets d'accès impose un niveau de sécurité maximal.

\subsection{Chiffrement des secrets et données sensibles}
Les identifiants de connexion aux systèmes cibles ne sont jamais stockés en clair. Nous utilisons l'algorithme \textbf{AES-256-GCM} pour le chiffrement symétrique des secrets. L'utilisation d'un \textit{Initialization Vector} (IV) unique pour chaque secret empêche les attaques par analyse de motifs et garantit la confidentialité des accès.

\subsection{Gestion des accès et authentification}
L'accès à l'interface d'administration est sécurisé par :
\begin{itemize}
    \item \textbf{JWT (JSON Web Tokens) :} Pour une authentification sans état, sécurisée et scalable entre le frontend Angular et le backend Spring Boot.
    \item \textbf{RBAC (Role-Based Access Control) :} Les permissions sont granulaires. Un profil \textit{Viewer} peut consulter le Hub 360, tandis qu'un profil \textit{Admin} a les droits de modifier les workflows et de déclencher des replays.
\end{itemize}

\section{Conclusion}
L'implémentation du Provisioning Hub a permis de transformer la vision théorique en une solution technique robuste. En combinant des patterns architecturaux modernes et des optimisations de bas niveau, nous avons construit un système capable de supporter des flux complexes avec une fiabilité industrielle. Le chapitre suivant sera consacré à la validation de ces résultats à travers une stratégie de tests rigoureuse.
 % Réalisation Technique et Implémentation
\chapter{Tests, Validation et Résultats}
\label{chap:tests}

Ce chapitre présente la démarche de validation adoptée pour s'assurer que le Provisioning Hub répond aux exigences fonctionnelles et techniques définies. Nous détaillerons la stratégie de tests, l'analyse des résultats et la présentation de l'interface utilisateur.

\section{Introduction}
La validation d'un système de synchronisation de données nécessite une approche multi-niveaux. Il ne s'agit pas seulement de vérifier que le code fonctionne, mais de garantir l'intégrité des données transférées, la stabilité du système sous charge et l'ergonomie de l'outil d'administration.

\section{Stratégie de Tests}
\label{sec:strategie_tests}

Nous avons mis en œuvre une pyramide de tests pour couvrir l'ensemble des risques techniques et fonctionnels.

\subsection{Tests Unitaires et d'Intégration}
L'utilisation de \textbf{JUnit 5} et \textbf{Mockito} a permis d'atteindre un taux de couverture élevé sur la logique métier. Un accent particulier a été mis sur le moteur de règles SpEL, avec des centaines de scénarios testant des expressions logiques complexes et des cas limites. Les tests d'intégration ont validé la chaîne complète : \textit{Réception Message $\rightarrow$ Traitement Domain $\rightarrow$ Persistance Postgres $\rightarrow$ Envoi Destination}.

\subsection{Tests de bout en bout (E2E) et Automatisation}
Pour valider l'expérience utilisateur et les flux transverses, nous avons utilisé \textbf{Selenium}. Des scripts automatisés simulent le parcours complet d'un administrateur :
\begin{enumerate}
    \item Création d'un nouveau workflow de provisionnement via l'interface Angular.
    \item Configuration d'une règle de filtrage SpEL complexe.
    \item Lancement d'une synchronisation manuelle.
    \item Vérification de l'apparition des logs en temps réel dans la vue Hub 360.
\end{enumerate}

\section{Analyse des Résultats}
\label{sec:analyse_resultats}

\subsection{Validation de la matrice de tests}
L'efficacité du système a été validée à travers une matrice de tests rigoureuse couvrant trois axes :
\begin{itemize}
    \item \textbf{Résolution d'Identité :} Le système a démontré une exactitude de 100\% dans le mapping des identifiants entre la source RH et les cibles LDAP/Jira, gérant parfaitement les cas de changements de noms ou de doublons.
    \item \textbf{Intégrité des Payloads :} Les données envoyées ont été validées par rapport aux schémas attendus par les API cibles. Aucun rejet pour format invalide n'a été observé après la phase de stabilisation.
    \item \textbf{Télémétrie et Audit :} La vue Hub 360 a permis de reconstruire l'historique complet de synchronisation pour n'importe quel collaborateur, validant l'objectif de traçabilité totale.
\end{itemize}

\subsection{Mesure des gains de performance et impact métier}
Le passage à l'approche déclarative a engendré des gains quantifiables et significatifs :
\begin{itemize}
    \item \textbf{Réduction du Time-to-Integration :} Le délai moyen pour ajouter une nouvelle destination est passé de 15 jours (cycle de dev complet) à moins d'une journée (configuration UI), représentant un gain d'agilité majeur.
    \item \textbf{Optimisation du débit :} Grâce au Bulk Loading, le temps de traitement d'un lot de 10 000 enregistrements a été divisé par quatre, permettant une synchronisation quasi instantanée des données.
\end{itemize}

\section{Présentation de l'Interface Administrateur}
\label{sec:interface}

L'interface, développée avec Angular 19, a été conçue pour simplifier la gestion d'un système complexe tout en offrant une visibilité maximale.

\subsection{Dashboard et Monitoring}
Le tableau de bord offre une vue synthétique de la santé du système. Il permet de surveiller en temps réel le nombre d'échanges en cours, le taux d'échec global et l'état des files d'attente RabbitMQ, permettant une réaction rapide en cas d'anomalie.

\subsection{Vue Employé 360 et Historique des échanges}
La fonctionnalité \textbf{Employee 360} est l'atout majeur du système. Elle agrège tous les \textit{Exchanges} liés à un individu. L'administrateur peut ainsi visualiser sur une seule page l'ensemble du cycle de vie du collaborateur : création dans l'AD, mise à jour dans Jira, et éventuels échecs de synchronisation avec le détail exact de l'erreur retournée par l'API.

\section{Conclusion}
La phase de validation confirme que le Provisioning Hub remplit l'ensemble des objectifs fixés initialement. La solution est robuste, performante et apporte une valeur ajoutée immédiate aux équipes opérationnelles en supprimant la dépendance au cycle de développement.
 % Tests, Validation et Résultats

% --- Conclusion et Annexes ---
\chapter*{Conclusion Générale et Perspectives}
\addcontentsline{toc}{chapter}{Conclusion Générale et Perspectives}

Ce projet a permis de concevoir et de réaliser le \textbf{Provisioning Hub}, une plateforme déclarative visant à moderniser la synchronisation des données RH vers divers systèmes cibles. En remplaçant un écosystème legacy fragmenté et rigide par une architecture orientée événements et un pipeline configurable, nous avons réussi à transformer un processus lourd et opaque en un flux agile, transparent et hautement résilient.

Le bilan technique et fonctionnel démontre que les objectifs ont été pleinement atteints. L'introduction du paradigme déclaratif a permis de réduire drastiquement le \textit{Time-to-Integration}, libérant les équipes de développement des tâches répétitives de configuration. L'implémentation de l'architecture hexagonale et du moteur de règles SpEL offre une flexibilité sans précédent, permettant l'adaptation du système aux évolutions métier en temps réel. Enfin, la mise en place du pattern \textit{Save-First} et de la vue \textit{Hub 360} a résolu le problème historique de la traçabilité, offrant une source unique de vérité pour l'audit et le diagnostic.

Cependant, tout système logiciel est évolutif. Ce travail ouvre la voie à plusieurs perspectives d'amélioration pour une version 2 :
\begin{itemize}
    \item \textbf{Conformité RGPD et Gouvernance des Données :} Intégration native du droit à l'oubli et mise en œuvre de mécanismes d'export automatique des données pour répondre aux exigences réglementaires européennes.
    \item \textbf{Transformations Avancées via JS Sandboxing :} Pour dépasser les limites de SpEL, l'intégration d'un moteur JavaScript isolé (sandbox) permettrait de réaliser des transformations de données extrêmement complexes tout en garantissant la sécurité du serveur.
    \item \textbf{Scalabilité Massive :} Optimisation du système pour supporter un volume dépassant les 165 000 enregistrements avec un temps de traitement minimal, notamment via l'implémentation du partitionnement de données (sharding) au niveau de la base de données.
\end{itemize}

En conclusion, le Provisioning Hub ne se limite pas à une simple amélioration technique, mais représente un véritable changement de paradigme dans la gestion du provisionnement RH. En alliant rigueur architecturale et flexibilité fonctionnelle, nous avons construit un outil capable de soutenir la croissance et la transformation digitale de l'organisation.

\chapter*{Références}
\addcontentsline{toc}{chapter}{Références}

\begin{itemize}
    \item \textbf{Documentation Spring Framework} : Spring Expression Language (SpEL) and Spring Boot 3.3.
    \item \textbf{Documentation Angular} : Angular 19 Framework and Reactive Forms.
    \item \textbf{PostgreSQL Documentation} : PostgreSQL 17 features and performance tuning.
    \item \textbf{RabbitMQ Documentation} : Advanced Message Queuing Protocol (AMQP) and DLQ implementation.
    \item \textbf{Project PRD} : Document de spécifications fonctionnelles et techniques du Provisioning Hub.
\end{itemize}

\chapter*{Annexes}
\addcontentsline{toc}{chapter}{Annexes}

\section{Extraits de code critiques}
\textit{(À compléter avec des snippets de code sur le Policy Engine ou le Save-First pattern)}

\section{Captures d'écran de l'interface utilisateur}
\textit{(À compléter avec les captures du Dashboard et de la vue Hub 360)}

\section{Configuration Docker Compose}
\textit{(À compléter avec le fichier docker-compose.yml)}

\section{Guide d'installation}
\textit{(À compléter avec les étapes de déploiement)}


\end{document}
